% 366_Feldstruktur_FFGFT_De_ch_kindle.tex
% Johann Pascher, 19. September 2026 -- Kindle 6x9-Fassung
% Felder in der FFGFT: Energie, Fluss und Geometrie

\providecommand{\BM}{\textbf{[B]}}
\providecommand{\KM}{\textbf{[K]}}
\providecommand{\SM}{\textbf{[S]}}
\providecommand{\SetzM}{\textbf{[SETZUNG]}}
\providecommand{\QM}{\textbf{[Q]}}
\providecommand{\EM}{\textbf{[E]}}
\providecommand{\XM}{\textbf{[X]}}

\chapter*{Felder in der FFGFT: Energie, Fluss und Geometrie}
\addcontentsline{toc}{chapter}{Felder in der FFGFT: Energie, Fluss und Geometrie}

\begin{abstract}
Dieses Dokument entwickelt einen einheitlichen Zugang zur Feldstruktur der FFGFT,
ausgehend von einem unerwarteten experimentellen Befund und endend bei den
Bahnresonanzen der Planeten. Der Einstieg ist absichtlich konkret: das Goubau-Leitungsexperiment
--- Energietransport auf einem einzigen Draht ohne Rückleiter --- zwingt dazu,
die Frage ernstzunehmen, wo Energie tatsächlich fließt.
Die Antwort ist in der klassischen Feldtheorie bereits vollständig vorhanden
und seit Poynting (1884) bekannt; sie wird in Einführungskursen aber regelmäßig
verdeckt durch das nützliche, aber vereinfachende Bild von Strom und Rückpfad.

Das Ziel ist nicht, neue Physik zu behaupten. Das Ziel ist, sichtbar zu machen,
dass Antennentechnik, Atomphysik, Kristallographie und Himmelsmechanik
strukturell dasselbe Prinzip verwirklichen: die Geometrie des Trägers
bestimmt die erlaubten Kopplungswinkel. Kopplung zwischen zwei Feldsystemen
ist nur dort möglich, wo ihre Modenstrukturen kompatibel sind. Inkompatible
Moden koppeln nicht --- unabhängig von der Stärke der Felder.

In der FFGFT ist dieser Zusammenhang nicht eine Analogie, sondern eine
Konsequenz der Grundformel $\tilde{T}\cdot m = 1$ auf dem kompakten Träger
$T^4/\mathbb{Z}_3$. Die Wicklungszahlen $(n_\theta, n_\varphi)$ auf diesem
Träger sind die topologischen Invarianten hinter den Auswahlregeln der
Atomphysik, den Bragg-Bedingungen der Kristallographie und den ganzzahligen
Bahnresonanzen der Himmelsmechanik. In FFGFT-Einheiten ($c=\hbar=\alpha=1$)
gilt $I=m$, $V=E$, $P=E\cdot m$ --- exakt, nicht als Näherung.

Die Antennentechnik dient dabei als Leitfaden, weil sie das am weitesten
ausgereifte Anwendungsfeld der elektromagnetischen Feldtheorie ist:
die Kopplungsbedingungen sind dort nicht abstrakt, sondern in Normen,
Messprotokollen und kommerziellen Simulatoren verankert. Was dort als
Friis-Gleichung, Impedanzanpassung und Arrayfaktor bekannt ist, erscheint
in der Atomphysik als Auswahlregel, im Kristall als Bragg-Bedingung
und im Sonnensystem als Laplace-Resonanz --- dieselbe Struktur, verschiedene Träger.
\end{abstract}

% ============================================================
\section*{Was dieses Dokument erklärt --- und was nicht}
\addcontentsline{toc}{section}{Was dieses Dokument erklärt --- und was nicht}

Vier Fragen stehen im Mittelpunkt, die in dieser Reihenfolge entwickelt werden:

\textbf{1. Wo fließt Energie?} Nicht im Leiter, sondern im Feld um ihn herum.
Das Goubau-Experiment erzwingt diese Antwort experimentell; der Poynting-Vektor
liefert sie mathematisch exakt. Strom und Spannung sind Projektionen, keine
primären Größen (§\,1--3).

\textbf{2. Was tut eine Antenne mit dem Feld?} Sie formt es aktiv um ---
im Sendefall durch aufgeprägte Randbedingungen, im Empfangsfall durch
selektiven Modenfilter. Passiv im Feld streut sie; nahe beieinander
deformieren sich zwei Antennen gegenseitig. Das ist keine Superposition,
sondern gegenseitige Verformung der Modenstrukturen (§\,4--9).

\textbf{3. Was ist ein Teilchen in diesem Bild?} Ein stabiler Feldknoten ---
eine lokalisierte Modenkonzentration auf dem Träger. Das Nahfeld des
Teilchens ist das Coulombfeld; das Fernfeld entsteht nur bei Beschleunigung
oder Zustandsübergang. Die Auswahlregeln der Atomphysik sind die
Mo\-den\-kom\-pa\-ti\-bi\-li\-täts\-be\-din\-gun\-gen dieses Knotens (§\,10).

\textbf{4. Gilt das Prinzip auf allen Skalen?} Ja --- von der Antenne über
das Atom und den Kristall bis zu den Planetenbahnen. Die Kopplungsbedingung
ist immer dieselbe: ganzzahlige Phasenkompatibilität der Modenspektren.
Die Skala ändert sich; die Struktur der Bedingung nicht (§\,11).

Was dieses Dokument \emph{nicht} erklärt: Warum der Träger $T^4/\mathbb{Z}_3$
ist (das ist \SetzM\ in Dok.~365); warum $\xi = 4/30\,000$ (das folgt aus
den Galois-Gruppenordnungen, Dok.~338/365); wie aus der Modenkompatibilität
konkrete Massen folgen (Dok.~314, 317--323). Diese Dokumente sind die
Quellen, auf die hier verwiesen wird. Dieses Dokument ist der Zugang vom
Konkreten zum Strukturellen.

\vspace{0.5em}
\noindent
\textit{Hinweis zur Lektüre:} Wer mit Antennentechnik vertraut ist, kann
§\,1--3 überfliegen und bei §\,4 einsteigen. Wer den FFGFT-Hintergrund kennt,
findet die Verbindung zur Modenstruktur in §\,5, §\,10 und §\,11.

% ============================================================
\section*{Zeichenerklärung}
\addcontentsline{toc}{section}{Zeichenerklärung}

\medskip
\textbf{Epistemische Marker:}
\begin{tabular}{@{}lp{4.5cm}@{}}
\textbf{[B]} & algebraisch bewiesen \\
\textbf{[K]} & numerisch bestätigt \\
\textbf{[S]} & strukturell plausibel \\
\textbf{[SETZUNG]} & motivierte Eingabe \\
\textbf{[E]} & etablierte Mathematik/Physik \\
\textbf{[X]} & geprüft, negativ \\
\end{tabular}

\medskip
\begin{longtable}{@{}lp{7cm}@{}}
\toprule
\textbf{Symbol} & \textbf{Bedeutung} \\
\midrule
\multicolumn{2}{l}{\textit{Felder und Energie}} \\
$F^{\mu\nu}$ & elektromagnetischer Feldstärketensor \\
$\mathbf{E}, \mathbf{B}$ & elektrisches und magnetisches Feld \\
$\mathbf{H}$ & magnetische Feldstärke ($\mathbf{H} = \mathbf{B}/\mu_0$ im Vakuum) \\
$\mathbf{S} = \mathbf{E}\times\mathbf{H}$ & Poynting-Vektor (Energieflussdichte, W/m$^2$) \\
$j^\nu$ & elektromagnetischer Vierstrom \\
$u$ & elektromagnetische Energiedichte \\
$\varepsilon_0,\,\mu_0$ & elektrische Feldkonstante, magnetische Feldkonstante \\
$\lambda$ & Wellenlänge ($\lambda = c/f$) \\
$\omega = 2\pi f$ & Kreisfrequenz \\
$k = 2\pi/\lambda$ & Wellenzahl \\
$\theta$ & Elevationswinkel (Richtdiagramm) oder Glanzwinkel (Bragg) \\
$\psi$ & Polarisationsfehlerwinkel zwischen Sender- und Empfangsantenne \\
\midrule
\multicolumn{2}{l}{\textit{FFGFT-Größen}} \\
$\tilde{T}$ & intrinsische Zeitskala; $\tilde{T} \cdot m = 1$ \\
$\xi = 4/30\,000$ & dimensionsloser Skalierungsparameter der FFGFT \\
$T^4/\mathbb{Z}_3$ & kompakter Orbifold-Träger der FFGFT \\
$(n_\theta, n_\varphi)$ & Wicklungszahlen auf $T^4$; topologische Invarianten \\
\midrule
\multicolumn{2}{l}{\textit{Antennentechnik}} \\
$R_\mathrm{str}$ & Strahlungswiderstand (Leistungsanteil in Abstrahlung) \\
$G$ & Antennengewinn (gegenüber isotropem Strahler) \\
$D$ & Direktivität (Gewinn ohne Verluste) \\
$A_\mathrm{eff}$ & effektive Apertur; $A_\mathrm{eff} = (\lambda^2/4\pi)\,G$ \\
$\eta_A$ & Aperturwirkungsgrad ($0 < \eta_A \leq 1$) \\
$G_S, G_E$ & Gewinn von Sender- bzw.\ Empfangsantenne (Friis-Gleichung) \\
$\Delta\theta_{3\,\mathrm{dB}}$ & Halbwertsbreite (\,3-dB-Keule) \\
$Z_\mathrm{ant}$ & Eingangsimpedanz der Antenne \\
$X_\mathrm{ant}$ & Blindwiderstand (imaginärer Teil von $Z_\mathrm{ant}$) \\
$\phi_n$ & Speisephase des $n$-ten Elements im Phased Array \\
$AF(\theta)$ & Arrayfaktor (Interferenzmuster des Phased Array) \\
$\varepsilon_r$ & relative Permittivität (Dielektrizitätszahl) \\
$\varepsilon_\mathrm{eff}$ & effektive Permittivität einer Mikrostreifenleitung \\
$\Delta L$ & Randkorrektur der Patchantennenlänge \\
$k_\mathrm{SPP}$ & Wellenzahl eines Oberflächenplasmon-Polaritons \\
$\varepsilon_m, \varepsilon_d$ & (komplexe) Permittivität von Metall bzw.\ Dielektrikum (SPP) \\
\midrule
\multicolumn{2}{l}{\textit{Atomphysik und Skalen}} \\
$\hbar$ & reduziertes Plancksches Wirkungsquantum \\
$\alpha \approx 1/137$ & Feinstrukturkonstante; in FFGFT-Einheiten $\alpha = 1$ \\
$n, l, m$ & Hauptquantenzahl, Drehimpulsquantenzahl, magnetische Quantenzahl \\
$E_n$ & Energieniveaus des Wasserstoffatoms; $E_n = -13{,}6\,\mathrm{eV}/n^2$ \\
$r_e \approx 2{,}82\,\mathrm{fm}$ & klassischer Elektronenradius \\
$\lambda_C = \hbar/m_e c$ & Compton-Wellenlänge des Elektrons (${\approx}386\,\mathrm{fm}$) \\
$a_0 \approx 52\,917\,\mathrm{fm}$ & Bohr-Radius \\
$\sigma_T$ & Thomson-Streuquerschnitt; $\sigma_T = \tfrac{8\pi}{3}r_e^2$ \\
$S_{AB}$ & Überlappintegral zweier Atomorbitale $\psi_A$, $\psi_B$ \\
$Q$ & Normalkoordinate einer Molekülschwingung \\
$\vec{\mu}$ & elektrisches Dipolmoment eines Moleküls \\
$d$ & Netzebenenabstand im Kristallgitter (Bragg-Bedingung) \\
$n_1, n_2, n_3$ & mittlere Bewegungen (Bahnfrequenzen) der Galileischen Monde \\
\bottomrule
\end{longtable}

% ============================================================
\section{Ausgangspunkt: Felder oder Strom?}
% ============================================================

In jedem Einführungskurs zur Elektrizitätslehre wird das Prinzip betont:
Strom braucht einen Rückpfad. Elektronen, die aus einem Leiter austreten,
müssen auf dem anderen zurückfließen. Dieses Bild ist für technische Zwecke
vollständig richtig --- es beschreibt Schaltkreise präzise. Es verbirgt aber
eine tiefere Frage: Wo fließt die Energie tatsächlich?

Das Goubau-Experiment \QM\ (G.~Goubau, 1950) demonstriert die Antwort
experimentell. Eine Hoch\-fre\-quenz\-si\-gnal\-quel\-le (${\sim}900\,\mathrm{MHz}$)
speist über einen kegelförmigen Übergang (Launcher) einen einzigen
isolierten Draht. Am anderen Ende wandelt ein identischer Konus (Catcher)
das Signal zurück in eine Koaxialleitung. Zwischen beiden Konen existiert
kein zweiter Leiter, kein Rückpfad im klassischen Sinn --- dennoch überträgt
die Anordnung das Signal mit geringen Verlusten \cite{Goubau1950}.

Eine physische Unterbrechung des Feldes (etwa eine Hand über dem Draht)
unterbricht die Übertragung sofort. Die Energie fließt im Feld um den Draht
herum, nicht im Draht selbst.

% ============================================================
\section{Der Poynting-Vektor als primäre Beschreibung}
% ============================================================

Die Energieflussdichte elektromagnetischer Felder ist der Poynting-Vektor
\cite{Stratton1941,Griffiths1999}:
\[
  \mathbf{S} = \mathbf{E} \times \mathbf{H}.
\]
Seine Richtung (Rechte-Hand-Regel: $\mathbf{E} \to \mathbf{H}$, Daumen in
Flussrichtung) zeigt an, wohin Energie transportiert wird. Die über den
Querschnitt eines Kabels integrierte Flussdichte ergibt exakt die
elektrische Leistung \EM:
\[
  P = \oint_A \mathbf{S} \cdot d\mathbf{A} = V \cdot I.
\]
Das ist keine Näherung, kein Zufall --- es ist eine exakte Identität \cite{Griffiths1999}.
Strom und Spannung sind makroskopische Projektionen des Feldgeschehens,
nicht die primären Träger der Energie.

\paragraph{Konsequenz.}
Für eine koaxiale Leitung verlaufen die $\mathbf{E}$-Feldlinien radial
vom Innenleiter zum Außenleiter; die $\mathbf{H}$-Feldlinien bilden
konzentrische Ringe. Das Kreuzprodukt $\mathbf{E}\times\mathbf{H}$ zeigt
axial --- entlang der Leitung. Die Energie fließt im Dielektrikum zwischen
Innen- und Außenleiter, nicht im Kupfer.

% ============================================================
\section{Die Goubau-Leitung: Felder ohne Rückleiter}
% ============================================================

Auf der Goubau-Leitung (G-line) gibt es keinen Außenleiter. Die
Feldkonfiguration unterscheidet sich dennoch nur graduell von der
Koaxialleitung, wenn man nahe genug am Draht misst \EM.

\subsection*{Feldkonfiguration}

Ein Spannungsimpuls, der sich auf dem Draht fortpflanzt, erzeugt
wechselnde positive und negative Potentialbereiche entlang des Leiters.
Die elektrischen Feldlinien verlaufen zwischen diesen Bereichen --- nicht
von Draht zu Rückleiter, sondern von einem Abschnitt des Drahtes zum
nächsten. Das Dielektrikum (die rote PVC-Isolierung) verlangsamt die
Ausbreitung der Felder und beugt die Feldlinien zur Drahtachse hin,
ähnlich wie Glas Lichtstrahlen bricht \cite{Goubau1950,Harms1907}. Dieser Mechanismus hält die
Felder in Drahtnähe --- er ist nicht kosmetisch, sondern physikalisch
notwendig für praktischen Betrieb.

\subsection*{Der Kegelübergang als Impedanzanpassung}

Die Kegel (Launcher / Catcher) passen die Wellenimpedanz der Koaxialleitung
($50\,\Omega$) an die Charakterimpedanz der G-line (typisch einige
$100\,\Omega$) an \EM. Im Inneren des Kegels dehnen sich die Feldlinien
vom koaxialen Muster auf das G-line-Muster aus. An der Öffnung des Kegels
schnappen die Feldlinien um und verbinden sich mit dem Draht. Der Rückpfad
für den Strom im Koaxialkabel wird durch die Verschiebungsströme im
Kegelinneren bereitgestellt --- die komplexen Felder im Kegel liefern
genau die richtigen Kräfte, um Elektronen in der äußeren Oberfläche zu
treiben.

\subsection*{Verluste}

Weil Strom im Draht fließt, treten $I^2 R$-Verluste auf --- genau wie
in einer konventionellen Zweidrahtleitung. Das zeigt: die G-line ist
kein reines Wellenleiterproblem. Der Energiefluss hat eine
Feldkomponente (Poynting) und eine Leitungskomponente (Ohm), beide
beschreiben dieselbe Physik aus verschiedenen Perspektiven.

% ============================================================
\section{FFGFT-Einheiten: Strom ist Masse, Spannung ist Energie}
% ============================================================

In FFGFT-Einheiten ($c = \hbar = \alpha = 1$, Dok.~011; Dok.~365 Schicht~2)
vereinfachen sich elektromagnetische Größen zu geometrischen Ausdrücken
in Masse und Zeit \KM:
\begin{align*}
  e &= 1 \quad \text{(Elementarladung dimensionslos)},\\
  V &= E/e = E \quad \text{(Spannung = Energie)},\\
  I &= e/T = 1/T = m \quad \text{(Strom = Masse, via }\tilde{T}\cdot m=1\text{)}.
\end{align*}
Die Leistungsgleichung lautet damit:
\[
  P = V \cdot I = E \cdot m,
\]
eine Beziehung, die direkt aus der Grundformel $\tilde{T}\cdot m=1$ folgt
(Dok.~365 Schicht~1) \KM.

\paragraph{Maxwell-Gleichungen in FFGFT-Einheiten.}
In FFGFT-Einheiten (Dok.~365 §\,2) nehmen die Maxwell-Gleichungen die Form an:
\[
  \partial_\mu F^{\mu\nu} = j^\nu, \qquad \partial_{[\mu}F_{\nu\rho]}=0.
\]
Keine Vorfaktoren, keine Kopplungskonstanten außer $j^\nu$ als geometrischem
Objekt. Die Kopplung ist in der Einheitenwahl absorbiert, nicht eliminiert;
der physikalische Inhalt ist identisch mit dem SI-System.

% ============================================================
\section{Feldlinien als geometrische Objekte auf $T^4/\mathbb{Z}_3$}
% ============================================================

In der FFGFT ist der physikalische Träger der kompakte Orbifold
$T^4/\mathbb{Z}_3$ (Dok.~365 Schicht~1) \SetzM. Felder sind Schnitte von
Vektorbündeln über diesem Träger; ihre Pe\-ri\-o\-di\-zi\-täts\-be\-din\-gun\-gen erzwingen
diskrete Modenspektren.

\subsection*{Feldlinien und Wicklungszahlen}

Auf einem Torus $T^4$ haben nicht-triviale Feldkonfigurationen
topologische Invarianten --- Wicklungszahlen $(n_\theta, n_\varphi)$,
die angeben, wie oft eine Feldlinie den Torus in einer Richtung umläuft
(Dok.~314; Dok.~247) \BM. Diese Wicklungszahlen klassifizieren die
Moden und bestimmen --- über $\tilde{T}\cdot m=1$ --- die Massenskala des
jeweiligen Sektors. Eine Feldlinie ist in dieser Beschreibung kein
Hilfskonstrukt, sondern ein topologisches Objekt mit physikalischer
Relevanz.

\subsection*{Feldlinien im G-line-Kontext}

Das Verhalten der Feldlinien in der Goubau-Leitung --- Ausdehnung im Kegel,
Umschnappen an der Öffnung, Anlagerung an den Draht --- ist aus dieser
Perspektive ein makroskopisches Beispiel für das Minimierungsprinzip des
Energieflusses. Bei jedem Orts- und Zeitpunkt wählt der Fluss den Weg,
der den Poynting-Vektor mit dem Träger konsistent hält. Das ist kein
mystisches Verhalten, sondern die kontinuierliche Version der topologischen
Optimierung auf $T^4/\mathbb{Z}_3$.

\subsection*{Feldform und Träger: keine Eins-zu-eins-Entsprechung}

Die Grundformel $\tilde{T}\cdot m=1$ macht keine Aussage über die
\emph{Form} eines Feldes im Raum. Sie legt fest, welche Massenskalen
zu welchen Zeitskalen gehören --- nicht, wie das Feld geometrisch
aussieht. Der kompakte Träger $T^4/\mathbb{Z}_3$ ist die Topologie,
die die erlaubten Randbedingungen vorgibt; die tatsächlichen
Feldkonfigurationen auf und um diesen Träger sind davon verschieden
und je nach Situation vielfältig verformt.

Ein Vergleich verdeutlicht dies: Die Länge einer Gitarrensaite bestimmt
das erlaubte Frequenzspektrum --- aber die Schallwelle im Umgebungsraum
ist nicht saitenförmig. Ebenso erzwingt $T^4/\mathbb{Z}_3$ das
Modenspektrum; was als Feld im dreidimensionalen Beobachtungsraum
erscheint, ist eine Projektion davon, die von den konkreten
Randbedingungen, der Umgebungsgeometrie und der Modenbesetzung
abhängt. Nahe am Leiter ist das Feld der G-line zylindrisch, weiter
weg elliptisch verformt, an Inhomogenitäten asymmetrisch --- nichts
davon ist torusförmig. Der Torus ist der interne geometrische Träger,
nicht die Gestalt des Feldes.

\subsection*{Wechselwirkung: Träger bestimmt erlaubte Kopplungswinkel}

Wenn zwei Felder unterschiedlicher Träger im Raum ineinandergreifen,
ist ihre Wechselwirkung nicht einfache Überlagerung. Jedes Feld trägt
seine eigene Modenstruktur mit sich, festgelegt durch die
Wicklungszahlen $(n_\theta, n_\varphi)$ des jeweiligen Trägers.
Kopplung ist nur dort möglich, wo die Phasenbedingungen beider
Modenspektren kompatibel sind --- wo die Wicklungszahlen ganzzahlig
zusammenpassen. Die Geometrie des Trägers bestimmt damit die
\emph{Winkel} der Wechselwirkung: unter welchen Phasenlagen
Energieübertrag stattfindet und unter welchen nicht \SM.

Das ist im Korpus bereits quantitativ ausgearbeitet: Die
Galois-Raster-Bedingung aus Dok.~328 (Kopplungsregime) und
Dok.~343~\S\,G$'$ legt fest, dass Kopplung nur bei Frequenzverhältnissen
im Galois-Raster $\{2,3,5,7,11,13\}$ effektiv ist; der Auflösungsboden
$100\xi$ ist die schärfste Grenze, unterhalb derer zwei Moden nicht
mehr als getrennt unterscheidbar sind. Die scheinbar freie Form
eines Feldes im Raum ist das Ergebnis dieser Kom\-pa\-ti\-bi\-li\-täts\-be\-din\-gun\-gen,
projiziert auf den Beobachtungsraum \SM.

% ============================================================
\section{Antennentechnik als ausgereifte Feldtheorie}
% ============================================================

Die Antennentechnik ist das am weitesten ausgereifte praktische Anwendungsgebiet
der elektromagnetischen Feldtheorie. Jahrzehnte an Messungen, Normen und
Berechnungsverfahren liegen vor; die Kopplungsbedingungen sind nicht abstrakt,
sondern in Handbüchern, in Messprotokollen und in kommerziellen Simulatoren
verankert. Genau deshalb ist sie der geeignetste Einstieg für das allgemeine
Prinzip: die Geometrie des Trägers bestimmt die erlaubten Kopplungswinkel.

\subsection*{Grundgrößen und ihre Feldinterpretation}

Jede Antenne ist ein Wandler zwischen geführter Welle (Koaxialkabel, Hohlleiter)
und freier Welle im Raum. Die entscheidenden Größen sind \EM:

\begin{itemize}[nosep]
\item \textbf{Strahlungswiderstand} $R_\mathrm{str}$: Der Widerstand, den die
  Antenne der geführten Welle präsentiert, weil sie Energie in den Raum abstrahlt.
  Für den Halbwellendipol: $R_\mathrm{str} \approx 73{,}1\,\Omega$.
  In FFGFT-Einheiten ($I=m$, $V=E$): $R_\mathrm{str} = P/I^2 = E/m^2 \cdot m = E/m$.
  Der Strahlungswiderstand ist damit eine Energie-pro-Masse-Relation --- ein
  direktes Echo von $\tilde{T}\cdot m=1$ \SM.

\item \textbf{Gewinn} $G$ und \textbf{Direktivität} $D$: Der Gewinn gibt an, um
  welchen Faktor die Strahlungsintensität in Hauptrichtung gegenüber einem
  isotropen Strahler erhöht ist. Er ist reines Feldgeometrieproblem: wie bündeln
  die Phasenbedingungen des Trägers (Antennenlänge, Abstände, Speisephase) den
  Poynting-Vektor in eine Vorzugsrichtung? Für den Halbwellendipol:
  $D = 1{,}64$ (2,15 dBi). Für eine Yagi-Uda-Antenne mit $n$ Elementen:
  $G \approx 10\log_{10}(n) - 1$ dBi als Faustformel.

\item \textbf{Effektive Apertur} $A_\mathrm{eff}$: Die Fläche, aus der eine
  Empfangsantenne Energie entnimmt. Zusammenhang mit Gewinn \EM:
  \[
    A_\mathrm{eff} = \frac{\lambda^2}{4\pi}\,G.
  \]
  Diese Formel ist reine Feldgeometrie: Wellenlänge und Gewinn bestimmen die
  effektive Wirkfläche, nicht die physische Größe der Antenne. Eine kleine
  Antenne mit hohem Gewinn hat dieselbe Apertur wie eine große mit niedrigem.
\end{itemize}

\subsection*{Friis-Übertragungsgleichung: Kopplung zweier Träger}

Die Friis-Gleichung beschreibt den Energieübertrag zwischen zwei Antennen
im Freifeld \cite{Friis1946,Balanis2016}:
\[
  \frac{P_E}{P_S} = G_S\,G_E\left(\frac{\lambda}{4\pi r}\right)^2
  \cos^2\psi,
\]
wobei $P_S$ die gesendete, $P_E$ die empfangene Leistung, $G_S, G_E$ die
Gewinne von Sender- und Empfangsantenne, $r$ der Abstand und $\psi$ der
Po\-la\-ri\-sa\-ti\-ons\-feh\-ler\-win\-kel sind.

Jeder Faktor dieser Gleichung entspricht einer Kom\-pa\-ti\-bi\-li\-täts\-be\-din\-gung
zwischen den Trägern:

\begin{center}
\begin{tabular}{@{}lp{2.8cm}p{2.8cm}@{}}
\toprule
\textbf{Faktor} & \textbf{Bedeutung} & \textbf{FFGFT} \\
\midrule
$G_S$ & Bündelung Sender & Modengeom.\ Quelle \\
$G_E$ & Bündelung Empfänger & Modengeom.\ Empf. \\
$(\lambda/4\pi r)^2$ & Freiraumdämpfung & Ausbreitungsfaktor \\
$\cos^2\psi$ & Polarisationsüberlapp & Wicklungskompatib. \\
\bottomrule
\end{tabular}
\end{center}

Der $\cos^2\psi$-Faktor ist dabei der direkteste Ausdruck des Prinzips:
bei orthogonaler Polarisation ($\psi=90°$) ist der Übertrag null ---
nicht wegen Absorption oder Verlust, sondern weil die Modengeometrien
unverträglich sind. Kein beliebig starker Sender überwindet das.
Das ist identisch mit der Aussage über Wicklungszahlen: wenn sie nicht
zusammenpassen, findet kein Energieübertrag statt \SM.

\subsection*{Impedanzanpassung als Phasenbedingung}

Damit eine Antenne maximal abstrahlt, muss ihre Eingangsimpedanz
konjugiert komplex zur Speiseimpedanz sein (Leistungsanpassung) \EM:
\[
  Z_\mathrm{Antenne} = R_\mathrm{str} + jX_\mathrm{Antenne}, \qquad
  Z_\mathrm{Quelle} = R_\mathrm{Quelle} - jX_\mathrm{Antenne}.
\]
Der Imaginärteil --- der Blindwiderstand --- ist ein Maß für gespeicherte
Feldenergie in der Nähe der Antenne (Nahfeld, reaktive Leistung). Ist er
nicht kompensiert, pendelt Energie zwischen Quelle und Antennennahfeld hin
und her, ohne abgestrahlt zu werden. Das ist dasselbe Phänomen wie die
reaktive Leistung auf der G-line (§\,3): Energie pulsiert im Nahfeld,
ohne Nettotransport.

Für die G-line war der Kegelübergang die Anpassungsstruktur (50\,$\Omega$
auf ${\sim}300\,\Omega$); für Antennen übernehmen das $\lambda/4$-Transformatoren,
Baluns oder Matching-Netzwerke. Der Mechanismus ist strukturell identisch:
Übergang zwischen zwei Modenspektren durch geometrisch erzwungene
Phasenbedingung.

\subsection*{Antennenarten als Modengeometrien}

Verschiedene Antennentypen realisieren verschiedene Modengeometrien.
Jeder Typ ist ein eigenständiger Träger mit einem charakteristischen
Modenspektrum; das Richtdiagramm im Fernfeld ist die Projektion dieses
Spektrums auf den Beobachtungsraum.

% -----------------------------------------------------------
\paragraph{1. Kurzdipolantenne ($\ell \ll \lambda$).}

Der Kurzdipol ist der einfachste Grenzfall: ein gerader Leiter, dessen
Länge weit unter der Wellenlänge liegt. Die Stromverteilung ist näherungsweise
dreieckig; das Fernfeld ergibt sich analytisch exakt \EM:
\[
  E_\theta \propto \frac{I_0\,\ell}{\lambda r}\,\sin\theta,
  \qquad
  F(\theta) = \sin^2\theta.
\]
Das Richtdiagramm $F(\theta) = \sin^2\theta$ ist eine Torus-Form um die
Dipolachse: maximale Abstrahlung senkrecht zur Achse ($\theta = 90°$),
null in Achsenrichtung ($\theta = 0°, 180°$). Direktivität $D = 1{,}5$
(1,76\,dBi). Strahlungswiderstand:
\[
  R_\mathrm{str} = 80\pi^2\left(\frac{\ell}{\lambda}\right)^2 \approx 790\,\Omega
  \cdot \left(\frac{\ell}{\lambda}\right)^2.
\]
Für $\ell = 0{,}1\lambda$: $R_\mathrm{str} \approx 7{,}9\,\Omega$ --- sehr
niedrig, daher starke Fehlanpassung an $50\,\Omega$-Systeme. Die
Modenstruktur ist einfach: eine einzige elektrische Dipolmode, deren
Wicklungsgeometrie die Torus-Form des Richtdiagramms erzwingt.

% -----------------------------------------------------------
\paragraph{2. Halbwellendipol ($\ell = \lambda/2$).}

Der Halbwellendipol ist der wichtigste Referenzstrahler der Antennentechnik.
Die Stromverteilung ist sinusförmig, das Richtdiagramm leicht schmaler als
beim Kurzdipol \EM:
\[
  F(\theta) = \frac{\cos^2\!\left(\frac{\pi}{2}\cos\theta\right)}{\sin^2\theta}.
\]
Direktivität $D = 1{,}64$ (2,15\,dBi). Strahlungswiderstand
$R_\mathrm{str} = 73{,}1\,\Omega$ --- nahezu ideal für $50\,\Omega$- oder
$75\,\Omega$-Koaxialkabel (Fehlanpassung $< 2{,}5:1$ ohne Anpassnetzwerk).
Die Halbwelle erzwingt durch ihre Resonanzbedingung ($\ell = \lambda/2$)
genau die Phasenlage, bei der Strahlungs- und Speisewiderstand ein reelles
Verhältnis eingehen: die imaginäre Komponente der Eingangsimpedanz
verschwindet bei Resonanz. Das ist eine ganzzahlige Mo\-den\-kom\-pa\-ti\-bi\-li\-täts\-be\-din\-gung
--- die Länge muss ein halbes ganzzahliges Vielfaches der Wellenlänge betragen.

% -----------------------------------------------------------
\paragraph{3. Monopolantenne ($\ell = \lambda/4$) über Massefläche.}

Ein Viertelwellen-Monopol über einer unendlich ausgedehnten leitenden Ebene
ist durch Spiegelung äquivalent zum Halbwellendipol: Die Massefläche liefert
das Spiegelbild. Das Richtdiagramm ist die untere Halbkugel des Dipol-Torus,
nach oben reflektiert; die gesamte Strahlungsleistung konzentriert sich auf
den oberen Halbraum \EM:
\[
  D_\mathrm{Monopol} = 2\,D_\mathrm{Dipol} = 3{,}28 \quad (5{,}15\,\mathrm{dBi}).
\]
Strahlungswiderstand $R_\mathrm{str} = 36{,}6\,\Omega$ (halb so groß wie beim
Dipol, da nur halber Raum). Praktische Ausführungen: Autoantenne, Mobilfunkmast \cite{Balanis2016}.
Der Gewinnzuwachs gegenüber dem Dipol entsteht nicht aus mehr Eingangsleistung,
sondern aus der geometrischen Spiegelungsbedingung --- Verdopplung des Raumwinkels,
der zur Abstrahlung beiträgt, ist eine reine Modengeometriefrage.

% -----------------------------------------------------------
\paragraph{4. Yagi-Uda-Antenne.}

Die Yagi-Uda-Antenne besteht aus einem gespeisten Dipol (Strahler), einem
etwas längeren passiven Reflektor und einer Reihe kürzerer passiver Direktoren
\cite{Yagi1928,Balanis2016}. Alle parasitären Elemente werden durch Induktion angeregt; ihre Länge
bestimmt die Phasenlage des induzierten Stroms:
\begin{itemize}[nosep]
  \item Reflektor ($\ell \approx 0{,}5\lambda + 5\,\%$): kapazitiv,
    re-radiiert mit Phasenverzug $\Rightarrow$ Auslöschung nach hinten.
  \item Direktoren ($\ell \approx 0{,}5\lambda - 5\,\%$): induktiv,
    re-radiieren mit Phasenvorsprung $\Rightarrow$ konstruktive Interferenz
    nach vorne.
\end{itemize}
Gewinn als Funktion der Direktoranzahl $n$ (Faustregel) \EM:
\[
  G \approx 10\log_{10}(0{,}8\,n + 1)\;\mathrm{dBd}
  \quad (n = \text{Anzahl Direktoren}).
\]
Beispielwerte: 3 Direktoren $\Rightarrow G \approx 7\,\mathrm{dBd}$,
10 Direktoren $\Rightarrow G \approx 10\,\mathrm{dBd}$.
Die Halbwertsbreite der Hauptkeule (in der $E$-Ebene) beträgt näherungsweise
\[
  \Delta\theta_{3\,\mathrm{dB}} \approx \frac{52°}{\sqrt{G_\mathrm{lin}}},
\]
wobei $G_\mathrm{lin}$ der lineare Gewinnfaktor ist. Das Richtdiagramm wird
also mit steigender Elementanzahl schmaler --- der Träger wird selektiver in
seinen Kopplungswinkeln.

% -----------------------------------------------------------
\paragraph{5. Hornantenne.}

Die Hornantenne ist der Übergang vom Hohlleiter in den freien Raum.
Der Hohlleiter trägt eine diskrete Modenstruktur (TE$_{10}$-Grundmode
bei Rechteckhohlleiter); der Horn erweitert diese Mode adiabatisch auf
eine ebene Welle im Fernfeld \cite{Balanis2016,Pozar2012}. Die Halbwertsbreite hängt direkt von
den Aperturabmessungen ab:
\[
  \Delta\theta_{E} \approx 56°\,\frac{\lambda}{a}, \qquad
  \Delta\theta_{H} \approx 67°\,\frac{\lambda}{b},
\]
wobei $a$ die $E$-Ebenen-Apertur und $b$ die $H$-Ebenen-Apertur sind.
Der Gewinn einer Pyramidalhorn-Antenne mit Apertur $A = a \times b$:
\[
  G = \eta_A\,\frac{4\pi A}{\lambda^2}, \qquad \eta_A \approx 0{,}5\text{--}0{,}7.
\]
Der Aperwirkungsgrad $\eta_A < 1$ entsteht durch die nicht-gleichförmige
Phasenverteilung in der Apertur (Kugelwelle statt ebene Welle am Hornmund).
Dies ist direkt die Modenunverträglichkeit zwischen Hohlleiter-Nahfeld und
Fernfeld-Ebenwelle: je besser die Phasenanpassung (flache Apertur), desto
höher $\eta_A$.

% -----------------------------------------------------------
\paragraph{6. Parabolantennen.}

Eine Parabolantenne mit Durchmesser $D$ und Brennweite $f$ bündelt einfallendes
Licht --- oder Mikrowellen --- im Fokus. Im Sendefall transformiert ein
Primärstrahler im Fokus eine Kugelwelle in eine ebene Wellenfront in der
Aperturebene. Das Fernfeld ist die Fouriertransformierte dieser Aperturverteilung
\EM. Für eine kreisförmige, gleichmäßig beleuchtete Apertur:
\[
  F(\theta) = \left[\frac{2\,J_1(\pi D\sin\theta/\lambda)}
                        {\pi D\sin\theta/\lambda}\right]^2,
\]
wobei $J_1$ die Bessel-Funktion erster Art ist \cite{Balanis2016}. Gewinn und Halbwertsbreite:
\[
  G = \eta_A\,\frac{\pi^2 D^2}{\lambda^2}, \qquad
  \Delta\theta_{3\,\mathrm{dB}} \approx 70°\,\frac{\lambda}{D}.
\]
Für eine 3-m-Parabolantenne bei 10\,GHz ($\lambda = 3\,\mathrm{cm}$):
$G \approx 10^4$ (40\,dBi), $\Delta\theta \approx 0{,}7°$. Die Bessel-Funktion
im Richtdiagramm ist keine Näherung, sondern die exakte Konsequenz der
Kreissymmetrie der Apertur: die Modengeometrie (Kreis) bestimmt vollständig
das Kopplungsmuster im Fernfeld.

% -----------------------------------------------------------
\paragraph{7. Patchantenne (Mikrostreifenleiter).}

Eine rechteckige Patchantenne der Länge $L$ und Breite $W$ über einer
Massefläche (Dielektrikum $\varepsilon_r$, Dicke $h$) resoniert, wenn
$L$ gleich einer halben effektiven Wellenlänge ist \cite{Pozar2012,Balanis2016}:
\[
  L = \frac{\lambda}{2\sqrt{\varepsilon_\mathrm{eff}}} - 2\Delta L,
\]
wobei $\varepsilon_\mathrm{eff}$ die effektive Dielektrizitätszahl und
$\Delta L$ die Randkorrektur ist:
\[
  \varepsilon_\mathrm{eff} = \frac{\varepsilon_r+1}{2}
    + \frac{\varepsilon_r-1}{2}\left(1 + \frac{12h}{W}\right)^{-1/2}.
\]
Das Richtdiagramm ist breitbandig in der $H$-Ebene (keine Bündelung), schmal
in der $E$-Ebene; typischer Gewinn 5--8\,dBi. Die erlaubten Moden sind
vollständig durch Geometrie und Material ($L$, $W$, $h$, $\varepsilon_r$)
bestimmt --- keine freien Parameter. Kopplung an eine externe Welle ist nur
möglich, wenn deren Polarisation und Einfallswinkel mit der Resonanzmode
des Patches kompatibel sind. Kreuz\-po\-la\-ri\-sa\-ti\-ons\-ent\-kopp\-lung ($> 30\,\mathrm{dB}$)
zeigt: orthogonale Polarisation koppelt praktisch nicht, auch wenn die
Frequenz exakt stimmt.

% -----------------------------------------------------------
\paragraph{8. Phased Array.}

Ein phased Array besteht aus $N$ identischen Strahlern mit individuell
steuerbarer Speisephase $\phi_n$. Das Gruppendiagramm ergibt sich als
Produkt aus Einzelstrahler-Richtdiagramm $F_\mathrm{elem}(\theta)$ und
Arrayfaktor \cite{Mailloux2017,Balanis2016}:
\[
  AF(\theta) = \sum_{n=0}^{N-1} \exp\!\left[j\left(n\,kd\cos\theta + \phi_n\right)\right],
\]
wobei $k = 2\pi/\lambda$ die Wellenzahl und $d$ der Elementabstand sind.
Durch Wahl von $\phi_n = -n\,kd\cos\theta_0$ schwenkt die Hauptkeule auf
den Winkel $\theta_0$ --- elektronische Strahlschwenkung ohne mechanische
Bewegung. Gewinn bei $N$ Elementen mit Abstand $d = \lambda/2$:
\[
  G_\mathrm{Array} = N \cdot G_\mathrm{elem}.
\]
Das Phased Array ist das reinste Beispiel für das Kopplungswinkelprinzip:
die Speisephasen $\phi_n$ sind ganzzahlige Vielfache eines Grundphasenschritts
--- sie bestimmen, unter welchem Winkel konstruktive Interferenz stattfindet.
Jede andere Winkelkombination ergibt destruktive Interferenz. In FFGFT-Sprache:
die $\phi_n$ sind diskrete Wicklungsphasenwerte; der Arrayfaktor selektiert
die kompatiblen Kopplungswinkel \SM.

% -----------------------------------------------------------
\paragraph{Übersichtstabelle.}

{\setlength{\hfuzz}{8pt}%
\begin{center}
\begin{tabular}{@{}lp{1.5cm}p{1.6cm}p{2.1cm}@{}}
\toprule
\textbf{Typ} & \textbf{Gewinn} & \textbf{$\Delta\theta$}
             & \textbf{Modengeom.} \\
\midrule
Kurzdipol     & 1,76\,dBi  & $90°$         & elektrische Dipolmode \\
Halbwellendipol & 2,15\,dBi & $78°$       & Resonanz $\ell=\lambda/2$ \\
Monopol ($\lambda/4$) & 5,15\,dBi & $46°$ & Spiegelmode über Masse \\
Yagi (5 Dir.) & ${\approx}12\,\mathrm{dBd}$ & ${\approx}30°$ & Interferenzarray \\
Horn & 15--25 & $3$--$15°$ & Apertur\\
Parab. 3m/10GHz & 40\,dBi & $0{,}7°$ & Fourier \\
Patch         & 5--8\,dBi  & $60°$--$120°$ & Resonanzmode im Dielektrikum \\
Phased ($N$) & $N G_\mathrm{e}$ & $\propto 1/\sqrt{N}$ & Phasenarray \\
\bottomrule
\end{tabular}
\end{center}}%

In jedem Fall gilt: die physische Geometrie der Antenne legt das erlaubte
Modenspektrum fest; was im Fernfeld als Richtdiagramm erscheint, ist die
Projektion dieses Spektrums auf den Beobachtungsraum. Kein Antennentyp
überträgt Energie unter Winkeln, die seiner Modengeometrie widersprechen ---
nicht näherungsweise, sondern strukturell. Das ist die technisch ausgereifte,
quantitativ verankerte Form des Prinzips, das in der FFGFT als Kompatibilität
der Wicklungszahlen formuliert ist \SM.

\subsection*{Reziproziät: Sender und Empfänger als duale Beschreibung}

Das Reziprozitätsprinzip der Antennentechnik besagt: Gewinn, Richtdiagramm
und Impedanz einer Antenne sind im Sende- und Empfangsfall identisch \cite{Balanis2016,Griffiths1999}.
Diese Symmetrie ist kein Zufall --- sie spiegelt die Zeitumkehrsymmetrie
der Maxwell-Gleichungen wider, und damit die Dualität von Feldquelle und
Feldsenke. In FFGFT-Sprache: Sender und Empfänger sind reziproke
Projektionen desselben Feldgeschehens, genau wie $\tilde{T}$ und $m$
reziproke Darstellungen desselben Objekts sind \SM.

% ============================================================
\section{Mehrere Feldträger im Raum: Randbedingungen, Kopplung und scheinbare Überlagerung}
% ============================================================

Das Superpositionsprinzip der klassischen Feldtheorie besagt: In einem linearen
Medium addieren sich Felder unabhängiger Quellen linear. Jedes Teilfeld breitet
sich aus, als wäre das andere nicht vorhanden. Das ist für weit voneinander
entfernte, schwach koppelnde Quellen eine ausgezeichnete Näherung --- und sie
funktioniert in der Praxis sehr gut. Sie verschleiert aber, was physikalisch
tatsächlich geschieht.

\subsection*{Das Gesamtfeld als konsistente Lösung aller Randbedingungen}

Jedes Objekt im Feld --- eine Antenne, ein Leiter, ein Dielektrikum, ein Atom
--- ist selbst ein Feldträger mit eigener Modenstruktur. Es erzwingt
Randbedingungen: die Tangentialkomponenten des elektrischen Feldes müssen
an jeder Leiteroberfläche verschwinden, die Normalkomponenten der
Flussdichte müssen an Grenzflächen stetig übergehen \EM. Das Gesamtfeld
ist nicht die Summe der Einzelfelder --- es ist das einzige konsistente
Feld, das alle Randbedingungen gleichzeitig erfüllt.

Die Superposition ist die Näherung, die entsteht wenn die Randbedingungen
der einzelnen Objekte räumlich weit genug getrennt sind, um sich nicht
gegenseitig zu beeinflussen. Sobald Objekte in Nahfeldabstand geraten,
bricht diese Näherung zusammen --- nicht graduell, sondern strukturell.

\subsection*{Nahfeld: der Ort wo Randbedingungen kollidieren}

Die Nahfeldzone einer Antenne reicht bis etwa $r < 2D^2/\lambda$
(Fraunhofer-Grenze), wobei $D$ die Antennengröße ist \EM. In dieser Zone
ist das Feld nicht als ebene Welle beschreibbar; es enthält reaktive
Anteile (gespeicherte Feldenergie, die hin- und herpulsiert, ohne
abgestrahlt zu werden). Bringt man eine zweite Antenne in diese Zone,
geschieht Folgendes:

\begin{itemize}[nosep]
\item Die Eingangsimpedanz beider Antennen ändert sich --- die
  Resonanzfrequenz verschiebt sich, der Strahlungswiderstand ändert
  sich. Das ist keine Addition von Feldern, sondern eine gegenseitige
  Deformation der Modenstrukturen.
\item Das Richtdiagramm jeder Antenne verändert sich, weil die
  Randbedingung der anderen Antenne die Feldverteilung im gemeinsamen
  Nahraum verformt.
\item Energie kann direkt zwischen den Nahfeldern übertragen werden
  --- ohne Abstrahlung in den Fernraum. Das ist das Prinzip der
  induktiven Kopplung (Transformator, Qi-Laden), der kapazitiven
  Kopplung und der resonanten Nahfeldkopplung (Kurs et al.\ 2007 \cite{Kurs2007},
  drahtloses Laden über Meter).
\end{itemize}

\subsection*{Mutual Coupling in Antennenarrays}

In einem Phased Array mit Elementabstand $d < \lambda$ ist die gegenseitige
Kopplung (mutual coupling) nicht vernachlässigbar \cite{Mailloux2017,Pozar2012}. Die Eingangsimpedanz
des $n$-ten Elements hängt von allen anderen ab:
\[
  Z_{nn}^{\mathrm{aktiv}} = Z_{nn}^{\mathrm{isoliert}}
  + \sum_{m \neq n} Z_{nm}\,\frac{I_m}{I_n},
\]
wobei $Z_{nm}$ die Gegeninduktanz zwischen Element $n$ und $m$ ist.
Das Gesamtverhalten des Arrays ist nur durch Lösung des vollständigen
Gleichungssystems beschreibbar --- nicht durch Addition von
Einzelelement-Feldern. In der Praxis wird mutual coupling durch
Optimierung der Elementgeometrie und Speisephasen kompensiert; es
ist ein Entwurfsproblem, das die Industrie mit vollständigen
elektromagnetischen Simulatoren (MoM, FEM, FDTD) löst.

\subsection*{Hohlraumresonator: das Gesamtfeld als Eigenmode}

Ein Metallhohlraum (Cavity) erlaubt nur diskrete Resonanzfrequenzen,
die durch seine Geometrie bestimmt sind \cite{Pozar2012,Griffiths1999}. Für einen Rechteckhohlraum
der Maße $a \times b \times c$:
\[
  f_{mnp} = \frac{c}{2}\sqrt{\left(\frac{m}{a}\right)^2
    + \left(\frac{n}{b}\right)^2 + \left(\frac{p}{c}\right)^2}.
\]
Bringt man ein Objekt in den Hohlraum --- einen Dielektrikum-Block,
einen Leiter, ein biologisches Gewebe --- verschieben sich alle
Resonanzfrequenzen. Das gesamte Modenspektrum des Hohlraums ändert
sich, nicht nur die Moden, die das Objekt direkt berühren. Das
Gesamtfeld ist die Eigenmode des zusammengesetzten Systems
(Hohlraum + Objekt), nicht die Summe der Eigenmode des leeren
Hohlraums und eines Streufeldes. Diese Verschiebung ist messbar
und wird in der Mikrowellen-Dielektrizitätsmessung genutzt
(Cavity-Perturbation-Methode).

\subsection*{Faraday-Käfig: Randbedingungen löschen ein Feld}

Der Faraday-Käfig ist der extremste Fall: eine geschlossene leitende
Hülle erzwingt, dass die Tangentialkomponente des elektrischen Feldes
auf ihrer Innenseite verschwindet. Das Außenfeld existiert unverändert
weiter; das Innenfeld wird zu null gezwungen \EM. Das Gesamtfeld ist
die eindeutige Lösung der Maxwell-Gleichungen mit dieser Randbedingung
--- keine Superposition, sondern eine Auslöschung durch erzwungene
Kompatibilität.

Unvollständige Schirmung (Lücken, endliche Leitfähigkeit) zeigt das
Prinzip in abgeschwächter Form: Felder dringen ein, aber nur mit den
Modengeometrien, die durch die Lückengeometrie erlaubt sind. Ein
Spalt der Breite $w$ lässt Felder mit $\lambda > 2w$ stark gedämpft
durch (Cut-off-Bedingung des Schlitzwellenleiters). Das ist wieder
Modenkompatibilität: nur Felder, deren Wellenlänge mit der
Spaltgeometrie kompatibel ist, koppeln durch.

\subsection*{Plasmonen: Feld und Materie untrennbar}

An der Grenzfläche zwischen Metall und Dielektrikum entstehen
Oberflächenplasmonen-Polaritonen (SPP) --- gekoppelte Moden aus
elektromagnetischem Feld und kollektiver Elektronenschwingung \cite{Maier2007}.
Die Dispersionsrelation lautet:
\[
  k_\mathrm{SPP} = \frac{\omega}{c}\sqrt{\frac{\varepsilon_m\,\varepsilon_d}
    {\varepsilon_m + \varepsilon_d}},
\]
wobei $\varepsilon_m$ die (komplexe) Permittivität des Metalls und
$\varepsilon_d$ die des Dielektrikums ist. Das SPP ist weder reines
elektromagnetisches Feld noch reine La\-dungs\-trä\-ger\-schwin\-gung --- es ist
eine gemeinsame Mode beider Träger. Eine Superposition wäre
bedeutungslos: das Feld existiert nur in Kopplung mit der Materie.

Das ist der schärfste verfügbare Beleg dafür, dass „mehrere Felder
im Raum" keine unabhängigen Objekte sind. Die Modenstruktur des
Gesamtsystems ist die physikalische Realität; die Zerlegung in
Teilfelder ist eine Rechenmethode, die unter bestimmten Bedingungen
funktioniert.

\subsection*{FFGFT-Deutung: ein Feld, viele Projektionen}

In der FFGFT gibt es auf dem Träger $T^4/\mathbb{Z}_3$ nur ein Feld
\SetzM. Was als „mehrere Felder" erscheint, sind Projektionen
verschiedener Modengruppen auf den Beobachtungsraum. Die scheinbare
Überlagerung ist eine Darstellungsfrage.

Wenn zwei Objekte mit verschiedenen Modenstrukturen in Wechselwirkung
treten, ist das Ergebnis nicht die Summe ihrer Felder, sondern die
gemeinsame Eigenmode des zusammengesetzten Trägers --- bestimmt durch
die Kompatibilität der Wicklungszahlen beider Träger. Das ist dasselbe
Prinzip, das in der Antennentechnik als mutual coupling, in der
Hohlraumphysik als Cavity-Perturbation und in der Nanophotonik als
Plasmon-Kopplung erscheint \SM.

Die Komplexität der Feldwechselwirkung ist damit keine Komplikation,
die man bei ausreichend großem Rechenaufwand loswerden kann --- sie ist
die direkte Konsequenz davon, dass jeder Träger dem Gesamtfeld
Randbedingungen aufzwingt. Je mehr Träger im Raum, desto mehr
gleichzeitig zu erfüllende Kom\-pa\-ti\-bi\-li\-täts\-be\-din\-gun\-gen, desto
komplexer das Gesamtfeld. Das ist nicht Superposition --- das ist
Geometrie.

% ============================================================
\section{Teilchen als Feldknoten: Nah- und Fernfeld an lokalisierten Feldstrukturen}
% ============================================================

In der klassischen Physik ist ein Teilchen ein Objekt, das ein Feld
\emph{besitzt} --- die Ladung erzeugt das Coulombfeld, die Masse erzeugt
das Gravitationsfeld. In der FFGFT ist diese Beschreibung umgekehrt: ein
Teilchen ist eine stabile, lokalisierte Feldkonfiguration --- ein Knoten
im globalen Feld auf dem Träger $T^4/\mathbb{Z}_3$ \SetzM. Die Masse ist
keine intrinsische Eigenschaft des Knotens, sondern folgt aus der
Grundformel $\tilde{T}\cdot m=1$: die Zeitskala der Feldkonfiguration
bestimmt die Masse (Dok.~365 Schicht~1, Dok.~124) \KM.

Diese Umkehrung hat direkte Konsequenzen für die Struktur des Feldes
um ein Teilchen --- und sie macht den Zusammenhang mit der Antennentechnik
sichtbar.

\subsection*{Stationäres Teilchen: reines Nahfeld}

Eine ruhende Punktladung $e$ erzeugt das Coulombfeld \EM:
\[
  \mathbf{E}(r) = \frac{e}{4\pi\varepsilon_0}\,\frac{\hat{r}}{r^2},
  \qquad
  \mathbf{B} = 0.
\]
In FFGFT-Einheiten ($e=1$, $4\pi\varepsilon_0=1$): $V(r) = -1/r$ \cite{PascherDok365}.
Dieses Feld fällt mit $1/r^2$ ab (Feldstärke) bzw.\ $1/r$ (Potential),
trägt keine Energie in den Fernraum und ist vollständig reaktiv ---
identisch zum Nahfeld einer stationären Antenne. Der Poynting-Vektor
$\mathbf{S} = \mathbf{E}\times\mathbf{H}$ ist für eine ruhende Ladung
null: kein Energiefluss, nur gespeicherte Feldenergie in der Knotenstruktur.

Das Coulombfeld ist das Nahfeld des Teilchens. Es verformt den Raum
um den Knoten herum, erzwingt Randbedingungen für alle anderen Felder
in seiner Umgebung --- genau wie eine passive Antenne das Feld in
ihrer Umgebung verformt, ohne selbst zu senden.

\subsection*{Beschleunigtes Teilchen: Nahfeld und abgestrahltes Fernfeld}

Sobald die Ladung beschleunigt wird, entsteht ein Fernfeld --- eine
sich ablösende elektromagnetische Welle. Die Larmor-Formel gibt die
abgestrahlte Leistung \cite{Jackson1999,Griffiths1999}:
\[
  P = \frac{e^2 a^2}{6\pi\varepsilon_0 c^3}
    = \frac{2}{3}\,\frac{e^2 a^2}{c^3}
  \quad \text{(in Gauß-Einheiten)},
\]
wobei $a$ die Beschleunigung ist. In FFGFT-Einheiten ($c=e=1$):
$P = \frac{2}{3}a^2$.

Das Fernfeld fällt mit $1/r$ ab (nicht $1/r^2$), transportiert Energie
dauerhaft weg und ist propagierend --- identisch zum Fernfeld einer
gespeisten Antenne. Die Grenze zwischen Nah- und Fernfeld liegt bei
$r \sim \lambda/2\pi$ (Übergangszone), wobei $\lambda = c/f$ und $f$
die charakteristische Frequenz der Beschleunigung ist.

Das ist das Antennenprinzip auf Teilchenebene:
\begin{itemize}[nosep]
  \item Stationäre Ladung = ungespeiste Antenne: nur Nahfeld,
    keine Abstrahlung.
  \item Gleichförmig bewegte Ladung = Antenne mit Gleichstrom:
    kein wechselnder Strom, keine Abstrahlung (Cherenkov-Strahlung
    ist ein Sonderfall für $v > c/n$).
  \item Beschleunigte Ladung = gespeiste Antenne: Nahfeld plus
    abgestrahltes Fernfeld, Leistung $\propto a^2$.
\end{itemize}

\subsection*{Atomarer Übergang: quantisiertes Fernfeld}

Im Atom sind die erlaubten Knotenstrukturen durch die Quantenzahlen
$(n, l, m)$ diskret --- das sind die erlaubten Moden des Coulombpotentials,
also die Eigenmoden des atomaren Trägers. Ein Übergang zwischen zwei
Zuständen ist eine Änderung der Knotenstruktur; dabei wird ein Photon
emittiert --- das abgestrahlte Fernfeld nimmt eine diskrete Energie
$\Delta E = \hbar\omega$ mit. Die Auswahlregeln \cite{CohenTannoudji1977,Griffiths1995}:
\[
  \Delta l = \pm 1, \qquad \Delta m = 0,\pm 1
\]
sind Kom\-pa\-ti\-bi\-li\-täts\-be\-din\-gun\-gen: nur Übergänge, bei denen die
Änderung der Modengeometrie mit der Geometrie eines Dipolphotons
kompatibel ist, sind erlaubt. Quadrupol- und höhere Übergänge
($\Delta l = \pm 2$) sind möglich, aber um Faktor $\alpha^2 \approx 10^{-4}$
unterdrückt --- weil die Modenkompatibilität schlechter ist.

Das ist in FFGFT-Sprache: die Wicklungszahlen des Anfangs- und
Endzustands müssen sich um genau die Wicklungszahlen eines Photons
unterscheiden. Die Auswahlregeln sind nicht postuliert, sondern
folgen aus der Geometrie des Trägers \SM.

\subsection*{Die ``Größe'' eines Teilchens: Ausdehnung des Knotens}

Ein Teilchen als Feldknoten hat keine scharfe geometrische Grenze.
Die ``Größe'' ist die räumliche Ausdehnung der Knotenstruktur,
also der Bereich, in dem das Nahfeld dominant ist gegenüber dem
Hintergrundfeld. Für das Elektron ist das der klassische
Elektronenradius \EM:
\[
  r_e = \frac{e^2}{4\pi\varepsilon_0 m_e c^2} \approx 2{,}82\,\mathrm{fm},
\]
der Compton-Wellenlänge $\lambda_C = \hbar/m_e c \approx 386\,\mathrm{fm}$
und dem Bohr-Radius $a_0 \approx 52\,917\,\mathrm{fm}$. Diese drei
Längen sind nicht willkürlich, sondern durch $\alpha$ verknüpft \cite{Jackson1999,PascherDok365}:
\[
  r_e = \alpha\,\lambda_C / 2\pi = \alpha^2 a_0.
\]
In FFGFT-Einheiten ($\alpha=1$) kollabieren diese drei Längen zu einer
einzigen --- die Hierarchie entsteht erst durch die SI-Einbettung via
$\xi$ (Dok.~365 Schicht~3) \KM.

Der Knoten ist also nicht punktförmig --- er hat eine Struktur, die
sich über mehrere Längenskalen erstreckt. Was als ``Punkt'' erscheint,
ist die Projektion dieser Struktur auf eine Messung, die grob genug
ist, um die innere Struktur nicht aufzulösen.

\subsection*{Vergleichstabelle: Antenne und Teilchen}

\begin{center}
\begin{tabular}{@{}lp{2.4cm}p{2.4cm}@{}}
\toprule
\textbf{Eigenschaft} & \textbf{Antenne} & \textbf{Teilchen} \\
\midrule
Nahfeld & reaktiv, $E\propto 1/r^2$--$1/r^3$ & Coulombfeld $E\propto 1/r^2$ \\
Fernfeld & abgestrahlte Welle, $E\propto 1/r$ & Photon bei Beschleunigung \\
Abstrahlung & nur bei Stromänderung & nur bei Beschleunigung/Übergang \\
Modenstruktur & Antennengeometrie & Quantenzahlen $(n,l,m)$ \\
Auswahlregel & Impedanzkompatibilität & $\Delta l=\pm1$, $\Delta m=0,\pm1$ \\
``Größe'' & physische Dimension & Knotenausdehnung $r_e,\lambda_C,a_0$ \\
Stationär & kein Fernfeld & kein Fernfeld (Coulomb only) \\
Passiv im Feld & Streuung, Abschattung & Polarisierbarkeit, Streuquerschnitt \\
\bottomrule
\end{tabular}
\end{center}

\subsection*{Streuquerschnitt als Maß der Knotenkopplung}

Ein Teilchen in einem externen Feld streut die einlaufende Welle ---
es verformt das Feld in seiner Umgebung, genau wie ein passiver Leiter.
Der Thomson-Streuquerschnitt des Elektrons \cite{Jackson1999}:
\[
  \sigma_T = \frac{8\pi}{3}r_e^2 \approx 6{,}65\times 10^{-29}\,\mathrm{m}^2
\]
ist die effektive Fläche, aus der das Elektron Energie aus dem
einlaufenden Feld entnimmt und reemittiert. Das ist direkt analog zur
effektiven Apertur $A_\mathrm{eff}$ einer Empfangsantenne (§\,8):
\[
  A_\mathrm{eff} = \frac{\lambda^2}{4\pi}G
  \quad \longleftrightarrow \quad
  \sigma_T = \frac{8\pi}{3}r_e^2.
\]
In beiden Fällen ist die Wirkfläche nicht die geometrische Größe des
Objekts, sondern das Ergebnis der Modenkompatibilität zwischen
einlaufendem Feld und Knotenstruktur. Für Photonen mit $\hbar\omega
\gg m_e c^2$ gilt die Compton-Streuung --- die Knotenstruktur des
Elektrons ist dann nicht mehr als punktförmig beschreibbar, und der
Querschnitt fällt mit steigender Energie.

\subsection*{FFGFT-Deutung: Knoten als lokale Modenkonzentration}

In der FFGFT ist ein Teilchen eine lokale Konzentration von Moden auf
dem Träger $T^4/\mathbb{Z}_3$ --- ein Ort, an dem bestimmte
Wicklungszahlen $(n_\theta, n_\varphi)$ besetzt sind und das Feld
lokal eine stabile, nicht-triviale Konfiguration einnimmt \SetzM.
Das Nahfeld des Teilchens ist die unmittelbare Umgebung dieses Knotens
im Träger, projiziert auf den Beobachtungsraum; das Fernfeld ist die
Abstrahlung, die entsteht wenn der Knoten seine Modenbesetzung ändert.

Die Grenze zwischen Nah- und Fernfeld ist damit keine scharfe Linie,
sondern der Übergang zwischen dem Bereich, in dem die Knotenstruktur
das Feld dominiert, und dem Bereich, in dem das Feld als freie Welle
propagiert. Dieser Übergang liegt bei $r\sim\lambda_C$ für relativistische
Beschreibungen und bei $r\sim a_0$ für atomare Übergänge --- beides
sind strukturelle Längenskalen des Knotens, keine freien Parameter \SM.

% ============================================================
\section{Skaleninvarianz der Kopplungsbedingung: Atome, Elemente und Planetenbahnen}
% ============================================================

Das Prinzip aus §\,8--10 --- Geometrie des Trägers bestimmt erlaubte
Kopplungswinkel, Kopplung nur bei Modenkompatibilität --- ist nicht auf
Antennen oder einzelne Teilchen beschränkt. Es gilt auf allen physikalischen
Skalen. Die Skala ändert sich; die Struktur der Bedingung bleibt identisch.
Dok.~242 (Skalenleiter) hat diesen Zusammenhang im Kosmologischen ausgeführt;
hier wird er von der Antennentechnik aus aufgespannt.

\subsection*{Atome: diskrete Moden im Coulombpotential}

Die Schrödinger-Gleichung im Coulombpotential hat diskrete Eigenmoden,
klassifiziert durch die Quantenzahlen $(n, l, m)$ \cite{Griffiths1995,CohenTannoudji1977}:
\[
  E_n = -\frac{m_e e^4}{2\hbar^2 n^2} = -\frac{13{,}6\,\mathrm{eV}}{n^2}.
\]
Das ist nichts anderes als die Resonanzbedingung eines Trägers: der
Träger ist das Coulombpotential (kugelsymmetrisch, $V\propto -1/r$);
die erlaubten Moden sind die Lösungen mit ganzzahligem $n$. Das Atom
ist ein Hohlraumresonator mit sphärischer Geometrie --- der Knoten
des Elektrons kann nur bestimmte stehende Wellenkonfigurationen
einnehmen, genau wie die TE$_{mn}$-Moden im Metallhohlraum (§\,9).

Die Auswahlregeln $\Delta l = \pm 1$, $\Delta m = 0,\pm 1$ (§\,10)
sind die Im\-pe\-danz\-an\-pas\-sungs\-be\-din\-gun\-gen dieses Resonators: ein
Dipolphoton ($l=1$) koppelt nur an Übergänge, bei denen die
Änderung der Knotengeometrie kompatibel mit einer Dipolmode ist.
Ein Quadrupolphoton ($l=2$) koppelt mit $\Delta l = \pm 2$, aber
schwächer um $\alpha^2$ --- weil die Modenüberlapp-Effizienz kleiner
ist, identisch zum Aperturwirkungsgrad $\eta_A < 1$ einer Hornantenne.

\subsection*{Periodensystem: Modenklassifikation eines Mehrteilchensystems}

Die Schalenstruktur des Periodensystems --- $1s^2\,2s^2\,2p^6\,3s^2\ldots$
--- ist eine Klassifikation nach Modengeometrien des effektiven
Zentralpotentials \EM:
\[
  \text{Schale } n: \quad 2n^2 \text{ Plätze}
  \quad (s, p, d, f \text{ mit Spin}).
\]
Jede Unterschale ($s, p, d, f$) hat eine andere Drehimpulsquantenzahl
$l$ und damit eine andere räumliche Modengeometrie. Die chemische
Bindung entsteht durch Kopplung der Randmoden zweier Atome: wenn die
äußersten besetzten Orbitale (HOMO/LUMO) eine kompatible Geometrie
haben --- ausreichender Raumüberlapp, kompatible Symmetrie --- entsteht
eine kovalente Bindung. Das Überlappintegral
\[
  S_{AB} = \int \psi_A^*(\mathbf{r})\,\psi_B(\mathbf{r})\,d^3r
\]
ist das direkte quantenmechanische Analogon des Polarisationsfaktors
$\cos^2\psi$ in der Friis-Gleichung (§\,8): $S_{AB}=0$ bedeutet
keine Kopplung, unabhängig von der Energie der Atome.

Die Edelgase (geschlossene Schalen, $S_{AB}\approx 0$ für alle
Bindungspartner) sind das quantenmechanische Analogon zur
orthogonalen Polarisation: kein Energieübertrag, weil keine
Modenkompatibilität.

\subsection*{Kristallgitter: Bragg-Bedingung als Kom\-pa\-ti\-bi\-li\-täts\-be\-din\-gung}

In einem Kristall bilden die Atome ein periodisches Gitter der
Gitterkonstante $a$. Eine elektromagnetische Welle (Röntgen, Neutronen,
Elektronen) wird nur dann kohärent reflektiert, wenn die
Bragg-Bedingung erfüllt ist \cite{Kittel2005,Ashcroft1976}:
\[
  2d\sin\theta = n\lambda, \quad n \in \mathbb{Z},
\]
wobei $d$ der Netzebenenabstand und $\theta$ der Glanzwinkel ist.
Das ist eine ganzzahlige Phasenbedingung: die Weglängendifferenz
zwischen an aufeinanderfolgenden Netzebenen reflektierten Teilwellen
muss ein ganzzahliges Vielfaches der Wellenlänge sein. Nur dann
addieren sich die Amplituden konstruktiv --- identisch zur
Phased-Array-Bedingung (§\,8):
\[
  \phi_n = -n\,kd\cos\theta_0.
\]
Das Kristallgitter ist ein natürliches Phased Array mit festem
Elementabstand $d$. Die Bragg-Reflexion selektiert genau die
Einfallswinkel, bei denen Modenkompatibilität herrscht; alle anderen
Winkel ergeben destruktive Interferenz. In der FFGFT entsprechen
die reziproken Gittervektoren $\mathbf{G}$ den erlaubten
Wicklungsvektoren des Trägers; Bragg-Reflexion ist die makroskopische
Manifestation der Quan\-ti\-sie\-rungs\-be\-din\-gung auf dem Torus \SM.

\subsection*{Molekülschwingungen: Infrarotaktivität als Auswahlregel}

Eine Molekülschwingungsmode ist nur dann infrarotaktiv (absorbiert
IR-Strahlung), wenn sie eine Änderung des elektrischen Dipolmoments
bewirkt \cite{Atkins2010,Wilson1955}:
\[
  \left(\frac{\partial \vec{\mu}}{\partial Q}\right)_{Q=0} \neq 0,
\]
wobei $Q$ die Normalkoordinate der Schwingung ist. Das ist die
Im\-pe\-danz\-an\-pas\-sungs\-be\-din\-gung zwischen Photon (Dipolstrahlung) und
Schwingungsmode: nur wenn die Schwingung eine Dipolgeometrie hat,
koppelt sie an ein Dipolphoton. Symmetrische Streckschwingungen (wie
in N$_2$ oder O$_2$) ändern das Dipolmoment nicht --- sie sind
IR-inaktiv, aber Raman-aktiv (koppeln an Quadrupolphotonen). Das
erklärt, warum N$_2$ und O$_2$ keine Treibhausgase sind: ihre
Schwingungsmoden sind modenunkompatibel mit IR-Dipolstrahlung.

\subsection*{Planetenbahnen: ganzzahlige Resonanzen im Gravitationsfeld}

Das Gravitationsfeld der Sonne ist der Träger des Planetensystems.
Seine Moden sind die Kepler-Bahnen, klassifiziert durch die
Bahnparameter $(a, e, i)$. Wenn zwei Planeten in der Nähe einer
ganzzahligen Resonanz $p:q$ ($p, q \in \mathbb{Z}$) sind, akkumulieren
sich kleine Gravitationsstörungen kohärent über viele Umläufe
(Dok.~242) \cite{PascherDok242,Murray1999}. Die Laplace-Resonanz der Galileischen Monde ist das
bekannteste Beispiel \cite{Murray1999,Peale1979}:
\[
  \frac{T_\mathrm{Io}}{T_\mathrm{Europa}} \approx \frac{1}{2},
  \qquad
  \frac{T_\mathrm{Europa}}{T_\mathrm{Ganymed}} \approx \frac{1}{2},
  \qquad
  n_\mathrm{Io} - 3n_\mathrm{Europa} + 2n_\mathrm{Ganymed} = 0
\]
(exakte Relation der mittleren Bewegungen $n = 2\pi/T$). Diese
Relation ist keine Koinzidenz: sie ist eine stabile Konfiguration,
weil die gegenseitigen Störungen bei exakter Resonanz keine
säkulare Energiezufuhr bewirken, sondern nur periodische. Das System
hat einen Attraktor bei der ganzzahligen Resonanz --- identisch
zum Attraktor einer Arnold-Zunge in der nichtlinearen Dynamik
(Dok.~328 §\,Kopplungsregime) \KM.

Weitere Beispiele:
\begin{itemize}[nosep]
  \item Jupiter/Saturn: $5:2$-Resonanz (``große Ungleichheit'',
    bekannt seit Laplace 1785) \QM.
  \item Neptun/Pluto: $3:2$-Resonanz (Plutinos) \QM.
  \item Kirkwood-Lücken im Asteroidengürtel: Lücken bei
    Resonanzen $1:3$, $1:4$ mit Jupiter (destruktive Resonanz,
    nicht stabilisierend) \QM.
  \item Saturnringe: Lücken (Cassini-Division) bei $2:1$-Resonanz
    mit Mimas \QM.
\end{itemize}

Die Kopplungsbedingung ist dieselbe wie in der Antennentechnik: ganzzahlige
Phasenkompatibilität. Der Träger ist verschieden (Gravitationsfeld statt
elektromagnetisches Feld), die Struktur der Bedingung ist identisch.

\subsection*{Übersichtstabelle: Kopplungsbedingung auf allen Skalen}

\begin{center}
\begin{tabular}{@{}lp{1.5cm}p{1.9cm}p{1.3cm}p{1.5cm}@{}}
\toprule
\textbf{Skala} & \textbf{Träger} & \textbf{Kopplung}
  & \textbf{Bsp.} & \textbf{Analogon} \\
\midrule
Antenne & Strahler & $\cos^2\psi$, $Z$ & Yagi & Friis \\
Atom & Coulomb & $\Delta l=\pm1$ & H-Lin. & Auswahlr. \\
Molekül & eff.~Pot. & $\partial\mu/\partial Q\neq0$ & IR & Dipol \\
Kristall & Gitter & $2d\sin\theta=n\lambda$ & Bragg & Array \\
Period. & Schalen & Überlapp $S_{AB}$ & Kovalen & Apertur \\
Mol. & Orbital & Symm. & HOMO & Mod. \\
Planeten & Gravitat. & $p\,n_1=q\,n_2$ & Laplace & Arnold \\
\bottomrule
\end{tabular}
\end{center}

\subsection*{Das gemeinsame Prinzip}

Alle Zeilen der Tabelle beschreiben dasselbe: ein System hat einen
Träger mit einem Modenspektrum; Kopplung an ein anderes System ist
nur dort möglich, wo die Modenspektren kompatibel sind; inkompatible
Moden koppeln nicht, unabhängig von der Stärke der beteiligten Felder.

Die technische Reife der Antennentechnik --- Jahrzehnte an Messungen,
normierten Verfahren, kommerziellen Simulatoren --- macht sie zur
greifbarsten Verwirklichung dieses Prinzips. Die Atommechanik,
die Festkörperphysik und die Himmelsmechanik sind dasselbe Prinzip
auf anderen Trägern. In der FFGFT ist der gemeinsame Träger
$T^4/\mathbb{Z}_3$; die Kopplungsbedingungen auf allen Skalen sind
Projektionen der Wicklungszahl-Kompatibilität auf den jeweiligen
Beobachtungsraum \SM.

% ============================================================
\section{Dualität: Feld und Teilchen, Energie und Strom}
% ============================================================

Die FFGFT behandelt Zeit-Masse-Dualität als fundamentales Prinzip
(Grundformel $\tilde{T}\cdot m=1$, Dok.~365) \SetzM. Dieselbe Dualität
erscheint in der elektromagnetischen Feldtheorie:

\begin{center}
\begin{tabular}{@{}lll@{}}
\toprule
\textbf{Feld-Sicht} & & \textbf{Strom-Sicht} \\
\midrule
Poynting-Vektor $\mathbf{S} = \mathbf{E}\times\mathbf{H}$ & $\longleftrightarrow$ & $P = V \cdot I$ \\
Feldenergie im Dielektrikum & $\longleftrightarrow$ & $I^2 R$ Joulescher Verlust \\
Evaneszente Welle außen & $\longleftrightarrow$ & Reaktive Leistung \\
Wicklungszahl (topologisch) & $\longleftrightarrow$ & Masse (dynamisch) \\
\bottomrule
\end{tabular}
\end{center}

In keinem dieser Paare ist eine Seite ``fundamentaler'' als die andere ---
sie beschreiben dasselbe Feldgeschehen aus verschiedenen Projektionen.
Das ist das präzise Analogon zur Zeit-Masse-Dualität in der FFGFT:
$\tilde{T}$ und $m$ sind reziproke Darstellungen desselben Objekts, nicht
zwei verschiedene Größen \KM.

\subsection*{Reaktive Leistung und evaneszente Felder}

Auf der G-line pulsiert die Energie in der Peripherie des Drahtfeldes
hin und her --- Energiefluss wechselt die Richtung im Takt des Feldes.
Das ist reaktive Leistung in der Wechselstromlehre, oder --- optisch
formuliert --- das Evaneszenzfeld auf der Schattensei\-te einer totalreflektierenden
Grenzfläche \EM. In der FFGFT entspricht dieses Verhalten dem nicht-propagierenden
Anteil der Fourier-Moden auf $T^4/\mathbb{Z}_3$: Moden unterhalb der
Spektrallücke (Dok.~314 §\,H) tragen keine Nettoenenergie, aber strukturieren
das Spektrum \SM.

% ============================================================
\section{Exakte Übereinstimmung: Poynting-Integral und Leistungsformel}
% ============================================================

Es ist lehrreich, die Exaktheit dieser Identität zu verstehen, weil sie in
Kursen oft als ``ungefähre Übereinstimmung'' oder ``Plausibilitätsargument''
präsentiert wird. Der Satz von Poynting zeigt \EM:
\[
  \frac{\partial u}{\partial t} + \nabla \cdot \mathbf{S} = -\mathbf{j}\cdot\mathbf{E},
\]
wobei $u = \frac{1}{2}(\varepsilon_0 E^2 + \mu_0^{-1}B^2)$ die
elektromagnetische Energiedichte ist. Integration über das Volumen des
Kabels mit den entsprechenden Randbedingungen liefert exakt $P = V \cdot I$
--- die einzige Voraussetzung ist die Gültigkeit der Maxwell-Gleichungen.

In FFGFT-Einheiten (ohne $\varepsilon_0, \mu_0, e$) reduziert sich die
Identität zu:
\[
  P = E \cdot m,
\]
was unmittelbar aus $V = E$, $I = m$ und der Grundformel $\tilde{T}\cdot m=1$
folgt \KM. Die Exaktheit ist strukturell, keine numerische Koinzidenz.

% ============================================================
\section{Zusammenfassung}
% ============================================================

\begin{enumerate}[nosep]
\item \textbf{Primäre Lokalisierung.} Energie fließt in den Feldern
  ($\mathbf{S} = \mathbf{E}\times\mathbf{H}$), nicht im Leiter.
  Spannung und Strom sind makroskopische Projektionen desselben
  Feldgeschehens \EM.

\item \textbf{Geometrische Objekte.} Feldlinien auf dem Träger
  $T^4/\mathbb{Z}_3$ sind topologische Objekte mit Wicklungszahlen;
  sie klassifizieren Moden und bestimmen Massenskalen über
  $\tilde{T}\cdot m=1$ \BM.

\item \textbf{FFGFT-Einheiten.} Mit $c=\hbar=\alpha=1$ gilt $I=m$,
  $V=E$, $P=E\cdot m$ --- exakt, nicht als Näherung.
  Die Maxwellgleichungen sind ohne Vorfaktoren \KM.

\item \textbf{G-line als Veranschaulichung.} Das Goubau-Experiment
  demonstriert makroskopisch, dass Energietransport ohne expliziten
  Rückleiter möglich ist, weil der Transport im Feld stattfindet.
  Die Kegel-Übergangsstruktur entspricht Impedanzanpassung;
  die Feldlinien optimieren den Energiefluss lokal \EM.

\item \textbf{Dualität.} Feld-Sicht und Strom-Sicht sind reziproke
  Projektionen --- kein Paar ist fundamentaler. Das spiegelt die
  Zeit-Masse-Dualität $\tilde{T}\cdot m=1$ auf der Ebene klassischer
  Felder \SM.
\end{enumerate}

\section*{Literatur}
\addcontentsline{toc}{section}{Literatur}

\begin{thebibliography}{99}

% --- Elektromagnetismus und Antennentechnik ---

\bibitem{Goubau1950}
G.~Goubau,
\textit{Surface waves and their application to transmission lines},
J.~Appl.~Phys.\ \textbf{21} (1950) 1119--1128.
Grundlegende Arbeit zur Einzel-Draht-Wellenleitung; Konus-Übergang,
Feldkonfiguration, Verlustmessung 6~dB/Meile.

\bibitem{Harms1907}
F.~Harms,
\textit{Elektromagnetische Wellen an einem Draht mit isolierender
zylindrischer Hülle},
Ann.~Phys.\ \textbf{23} (1907) 44--60.
Erste Analyse des Dielektrikums als feldführende Schicht; Vorläufer der
Goubau-Leitung.

\bibitem{Stratton1941}
J.~A.~Stratton,
\textit{Electromagnetic Theory},
McGraw-Hill, New York, 1941.
Standardreferenz für Poynting-Satz (Kap.~2), Randbedingungen, Feldenergie.

\bibitem{Griffiths1999}
D.~J.~Griffiths,
\textit{Introduction to Electrodynamics},
3.~Aufl., Prentice Hall, 1999.
Poynting-Vektor, Larmor-Formel, Strahlungswiderstand, Reziprozität;
zugängliche Herleitung aller in §\,2--5 verwendeten Standardresultate.

\bibitem{Jackson1999}
J.~D.~Jackson,
\textit{Classical Electrodynamics},
3.~Aufl., Wiley, 1999.
Larmor-Formel (Kap.~14), Thomson-Streuquerschnitt (Kap.~14.8),
klassischer Elektronenradius, Compton-Wellenlänge; kanonische Referenz.

\bibitem{Balanis2016}
C.~A.~Balanis,
\textit{Antenna Theory: Analysis and Design},
4.~Aufl., Wiley, 2016.
Standardwerk der Antennentechnik: Richtdiagramme, Gewinn, Strahlungswiderstand,
Kurzdipol, Halbwellendipol, Monopol, Yagi-Uda, Hornantenne, Paraboloid,
Patchantenne, Phased Array, Friis-Gleichung, Reziprozitätssatz.

\bibitem{Friis1946}
H.~T.~Friis,
\textit{A note on a simple transmission formula},
Proc.~IRE \textbf{34} (1946) 254--256.
Originalarbeit zur Übertragungsgleichung;
``\textit{The received power is [...] proportional to the product of the
transmitting and receiving gains.}''

\bibitem{Yagi1928}
H.~Yagi,
\textit{Beam transmission of ultra short waves},
Proc.~IRE \textbf{16} (1928) 715--741.
Originalarbeit zur Yagi-Uda-Antenne; Reflektor- und Direktor-Prinzip;
Gewinn durch Phaseninterferenz passiver Elemente.

\bibitem{Pozar2012}
D.~M.~Pozar,
\textit{Microwave Engineering},
4.~Aufl., Wiley, 2012.
Hohlraumresonator (Kap.~6), Patchantenne (Kap.~14), Hornantenne,
Hohlleitermoden, Impedanzanpassung, Smith-Diagramm, Cavity-Perturbation.

\bibitem{Mailloux2017}
R.~J.~Mailloux,
\textit{Phased Array Antenna Handbook},
3.~Aufl., Artech House, 2017.
Arrayfaktor, Strahlschwenkung, mutual coupling, aktive Impedanzmatrix,
Grating-Lobes-Bedingung; quantitative Behandlung aller Array-Phänomene.

\bibitem{Maier2007}
S.~A.~Maier,
\textit{Plasmonics: Fundamentals and Applications},
Springer, 2007.
Oberflächenplasmonen-Polaritonen, Dispersionsrelation (Kap.~2),
Kopplungsbedingung, Nahfeldverstärkung; Standardwerk der Nanophotonik.

\bibitem{Kurs2007}
A.~Kurs, A.~Karalis, R.~Moffatt, J.~D.~Joannopoulos,
P.~Fisher, M.~Soljačić,
\textit{Wireless power transfer via strongly coupled magnetic resonances},
Science \textbf{317} (2007) 83--86.
Resonante Nahfeldkopplung über 2~m mit 40~\% Effizienz;
``\textit{The transfer is mediated by the overlap of the near fields.}''
Direkter Beleg für Energieübertrag ohne Abstrahlung durch Modenkompatibilität.

% --- Quantenmechanik und Atomphysik ---

\bibitem{Griffiths1995}
D.~J.~Griffiths,
\textit{Introduction to Quantum Mechanics},
Prentice Hall, 1995.
Wasserstoffatom, Quantenzahlen, Auswahlregeln $\Delta l=\pm1$,
$\Delta m=0,\pm1$ (Kap.~9); Überlappintegral und Bindungstheorie (Kap.~7).

\bibitem{CohenTannoudji1977}
C.~Cohen-Tannoudji, B.~Diu, F.~Laloë,
\textit{Quantum Mechanics}, 2~Bde.,
Wiley, 1977.
Auswahlregeln (Kap.~XIII), Dipolnäherung, Quadrupolunterdrückung
$\sim\alpha^2$; kanonische Referenz für atomare Übergänge.

% --- Festkörperphysik ---

\bibitem{Kittel2005}
C.~Kittel,
\textit{Introduction to Solid State Physics},
8.~Aufl., Wiley, 2005.
Bragg-Bedingung (Kap.~2), reziproker Gitterraum, Brillouin-Zone,
Gitterdynamik; Standardwerk der Festkörperphysik.

\bibitem{Ashcroft1976}
N.~W.~Ashcroft, N.~D.~Mermin,
\textit{Solid State Physics},
Holt, Rinehart and Winston, 1976.
Bragg-Streuung, Phased-Array-Analogie des reziproken Gitters (Kap.~6),
Elektronenbänder; fortgeschrittene Referenz.

% --- Molekülphysik ---

\bibitem{Atkins2010}
P.~Atkins, J.~de~Paula,
\textit{Physical Chemistry},
9.~Aufl., Oxford University Press, 2010.
IR-Aktivitätsbedingung $\partial\vec{\mu}/\partial Q\neq0$ (Kap.~13),
Raman-Aktivität, Symmetrieauswahlregeln; Standardwerk der physikalischen Chemie.

\bibitem{Wilson1955}
E.~B.~Wilson, J.~C.~Decius, P.~C.~Cross,
\textit{Molecular Vibrations},
McGraw-Hill, 1955.
Normalkoordinaten, Auswahlregeln, Gruppentheorie der Schwingungsmoden;
klassische Originalreferenz.

% --- Himmelsmechanik ---

\bibitem{Murray1999}
C.~D.~Murray, S.~F.~Dermott,
\textit{Solar System Dynamics},
Cambridge University Press, 1999.
Bahnresonanzen (Kap.~8), Laplace-Resonanz Io/Europa/Ganymed,
Kirkwood-Lücken, Cassini-Division, Arnold-Zungen in der
Himmelsmechanik; maßgebliche moderne Referenz.

\bibitem{Peale1979}
S.~J.~Peale, P.~Cassen, R.~T.~Reynolds,
\textit{Melting of Io by tidal dissipation},
Science \textbf{203} (1979) 892--894.
Laplace-Resonanz als Ursache der Tidal-Heizung von Io;
``\textit{The forced eccentricity [...] results from the Laplace relation
$n_1 - 3n_2 + 2n_3 = 0$.}'' Spektakuläre Vorhersage vor Voyager-1-Ankunft.

% --- FFGFT-Korpus (Interlinks) ---

\bibitem{PascherDok365}
J.~Pascher,
\textit{Dok.~365: Das Fundament der FFGFT in vier Schichten ---
Grundformel, elektromagnetischer Sektor, Skalierungsparameter und
algebraische Bestätigungen},
FFGFT-Korpus, Zenodo (2026).
Schicht~1 ($\tilde{T}\cdot m=1$), Schicht~2 (FFGFT-Einheiten,
$I=m$, $V=E$, Maxwell ohne Vorfaktoren), Schicht~3 ($\xi=4/30000$).

\bibitem{PascherDok314}
J.~Pascher,
\textit{Dok.~314: Gitter im Hilbertraum},
FFGFT-Korpus, Zenodo (2026).
Fourier-Moden auf $T^4/\mathbb{Z}_3$, Wicklungszahlen als topologische
Invarianten, Schalenzählung bis Norm~120, Spektrallücke.

\bibitem{PascherDok328}
J.~Pascher,
\textit{Dok.~328: Kopplungsregime, Resonanz und Synchronisation},
FFGFT-Korpus, Zenodo (2026).
Galois-Raster-Bedingung, Arnold-Zungen, Adler-Modell, $K_\mathrm{eff}=2\pi\xi$,
Kopplungsschwelle; direkte Verbindung zu Bahnresonanzen und mutual coupling.

\bibitem{PascherDok242}
J.~Pascher,
\textit{Dok.~242: Skalenleiter --- interskalare Kopplungen in der FFGFT},
FFGFT-Korpus, Zenodo (2026).
Planetenbahn-Argument, Laplace-Resonanz als Strukturprinzip,
Skalentrennung als historische Konvention, nicht physikalische Notwendigkeit.

\bibitem{PascherDok124}
J.~Pascher,
\textit{Dok.~124: Ladung als geometrische Projektion},
FFGFT-Korpus, Zenodo (2026).
Ladung als emergente Projektion geometrischer Muster; Teilchen als
Feldkonfiguration, nicht als intrinsische Punkteigenschaft.

\bibitem{PascherDok343}
J.~Pascher,
\textit{Dok.~343: Die Zeta-Funktion des $T^4/\mathbb{Z}_3$-Torus},
FFGFT-Korpus, Zenodo (2026).
Auflösungsboden $100\xi$, Galois-Raster, Farey-kritische Nennerhöhe;
quantitative Grenze der Mo\-den\-un\-ter\-scheid\-bar\-keit.

\end{thebibliography}
